\begin{recipe}{Ravioli met ricotta en spinazie}
\kicker[Hoofdgerecht,Vegetarisch,Italiaans]{Hoofdgerecht · vegetarisch · Italiaans}
\index[register]{Ravioli met ricotta en spinazie}
\index[register]{Ravioli}
\index[register]{Ricotta}
\index[register]{Spinazie}

\meta{100 minuten \dvd Voor 4 personen \dvd Raviolimaker}

\lettrine{R}{icotta} en spinazie zijn een klassieke combinatie voor gevulde
pasta. Het deeg maak je zelf van ei en bloem, de vulling is eenvoudig en mild.
Het meeste werk zit in het dun uitrollen en vullen. Met de raviolimaker maak je ze twaalf tegelijk.

\blockrule{Ingrediënten}
\begin{ingredients}
  \ing{70 g}{ricotta}\index[register]{Ricotta}
  \ing{135 g}{verse spinazie}\index[register]{Spinazie}
  \ing{2}{eieren}\index[register]{Eieren}
  \ing{200 g}{bloem tipo 00}\index[register]{Bloem}
  \ing{30 g}{Parmezaanse kaas, geraspt}\index[register]{Parmezaanse kaas}
  \ingb{olijfolie}
  \ingb{griesmeel, om te bestuiven}
  \ingb{water}
  \ingb{nootmuskaat, vers geraspt}
  \ingb{zout}
  \ingb{versgemalen zwarte peper}
\end{ingredients}

\blockrule{Bereiding}
\tip{Knijp de spinazie goed uit, anders maakt het overtollige vocht de ravioli
open tijdens het koken. Zelf ricotta maken kan ook, zie het basisrecept op
p.~\pageref{recipe:ricotta}. Lekker met gesmolten boter en salie, of met een
simpele tomatensaus.}
\begin{steps}
  \item Klop de eieren los in een kom. Doe de bloem in een grote kom of op een
        schone werkbank en maak een kuiltje. Giet de eieren er beetje bij beetje
        in en meng met een vork tot er een grof deeg ontstaat.
  \item Kneed het deeg een paar minuten tot het glad en gelijkmatig is.
        Wikkel het in vershoudfolie en laat 30 minuten rusten op een koele,
        droge plek.
  \item Was de spinazie. Verhit een scheutje olijfolie in een pan op
        middelhoog vuur en bak de spinazie 1 minuut tot hij begint te slinken.
        Zet het vuur laag, leg het deksel erop en laat 5 tot 6 minuten zachtjes
        garen. Breng op smaak met zout.
  \item Doe de spinazie in een vergiet en druk hem goed uit tot er geen vocht
        meer uitkomt. Hak hem daarna fijn.
  \item Meng in een kom de ricotta en de Parmezaanse kaas. Breng op smaak met
        peper en nootmuskaat en roer de spinazie erdoor.
  \item Halveer het deeg en houd de ene helft goed afgedekt, zodat die niet
        uitdroogt. Bestuif de andere helft licht met griesmeel, druk hem plat en
        rol hem met de deegroller uit tot een rechthoek zo breed als de
        pastamachine.
  \item Draai de pastamachine op de breedste stand en haal het deeg erdoor. Vouw
        de zijkanten naar het midden, rol de baan met de deegroller weer plat
        en haal hem opnieuw door de machine. Herhaal dit en zet de machine
        steeds dunner tot het deeg 1 tot 2 mm dik is. Wordt het deeg te nat,
        bestuif het dan met wat griesmeel. Rol de andere helft op dezelfde
        manier uit en dek de eerste baan intussen af met een theedoek.
  \item Snijd de banen met een mes in stukken die iets breder en langer zijn
        dan de raviolimaker. De restjes deeg kun je opnieuw kneden en
        uitrollen, als ze maar niet uitdrogen.
  \item Bestuif de raviolimaker met bloem en leg een deegbaan erop. Druk het
        deeg met de kunststof vorm in de vakjes. Schep in elk vakje ongeveer
        een theelepel vulling.
  \item Bevochtig het deeg rond de vulling met een natte vinger en leg de
        tweede deegbaan erover. Rol met de bijgeleverde deegroller met lichte
        druk over de maker, zodat de ravioli aan de randen worden dichtgedrukt
        en afgesneden.
  \item Keer de raviolimaker om boven een met griesmeel bestrooide theedoek, de
        twaalf ravioli komen er zo uit. Herhaal dit met de rest van het deeg
        en de vulling.
  \item Kook de ravioli in ruim kokend gezouten water op middelhoog vuur. Ze
        zijn klaar zodra ze bovendrijven, na 2 tot 3 minuten. Schep ze met een
        schuimspaan uit het water.
\end{steps}
\end{recipe}
